\section{问题三：通信约束下的运输与中继联合调度}\label{sec:q3}

\subsection{模型名称与新增约束}

在问题二的调度之上引入\textbf{连续通信约束}：运输无人机在\textbf{爬升、巡航、下降及物资投送}四个阶段均应保持通信。形式化地，设架次 $p$ 的轨迹为 $\mathrm{Traj}_{p}(t)$，则要求
\begin{equation}
	\forall t\in[t_{start}(p),\,t_{end}(p)]:\quad
	\mathrm{comm\_state}\big(\mathrm{Traj}_{p}(t)\big)\neq\text{中断}
	\label{eq:contcomm}
\end{equation}
其中通信状态按式~\eqref{eq:tristate} 的三态规则判定。该约束带来两个新难点：\textbf{①判定必须逐时刻}——只检查服务区一点或航段两端点会漏掉“中途被山体遮挡”的时段；\textbf{②选址是连续三维变量}——悬停位置含平面坐标与飞行海拔，若逐时刻独立选址，问题退化为难解的连续最优控制。

\subsection{新增决策变量}

空中中继的通信能力与悬停高度直接相关，已有研究给出了最优悬停高度与覆盖半径的解析关系\cite{alhourani2014lap}；无人机通信的总体机遇与挑战见文献\cite{zeng2016uavcomm}。

中继无人机的\textbf{悬停位置}（平面位置 $+$ 飞行海拔）、\textbf{服务时段}、\textbf{能源组件分配}。约束包括：悬停位置须位于 DEM 覆盖范围内；悬停\textbf{离地高度不超过 \SI{300}{m}}；往返悬停位置的地形净空按航段规则确定；须满足返航电量要求。

\subsection{求解方法：把“连续轨迹 $\times$ 连续选址”降维}

本文采用一个\textbf{可验证的充分条件}替代连续最优控制：对某架次，若存在\textbf{某一个}悬停点 $P$，使得轨迹上每一时刻都满足
\begin{equation}
	\text{直连可用}\ \vee\ \big(\text{接入}(P)\ \text{可用}
	\ \wedge\ \text{回传}(P)\ \text{可用}\big)
	\label{eq:suffcond}
\end{equation}
则该架次由一架中继在 $P$ 处全程保障即可。这样只需“离散候选点 $\times$ 时间采样”两次扫描。若单点无法覆盖全程，则退化为\textbf{按时段分段接力}——题目允许不同时刻由不同中继保障，只要每一时刻仅由一架中继服务。

\textbf{实现要点}：在 DEM 覆盖范围内、服务区凸包外扩 \SI{3}{km} 的区域上以 \textbf{\SI{400}{m} 网格}生成 \textbf{1885 个}悬停候选点；对每个需保障架次，在能覆盖其全部中断样本且返航 SOC 达标的候选点中选能耗最低者。为避免在 1885 个候选点上做全轨迹扫描，算法先用一次扫描筛出直连中断样本（约占 0\%$\sim$56.2\%），选址只在中断样本上评估；覆盖判定用稀疏探针快速否决。采样步长取 $\Delta t=\SI{2}{s}$。

\textbf{服务窗口的确定（易错点）}：需要中继驻留的时段应取该架次的\textbf{中断区间集合}，而\textbf{不是整条轨迹、也不是“首个中断样本到末个中断样本”的包络}。原因是一个架次的中断可能分成好几段（实测 2$\sim$10 段），中间夹着\textbf{直连本身可用}的时段，那些时段并不需要中继。实测包络会把站岗时长\textbf{虚高 39.3\%}（\SI{277.4}{min} 对 \SI{168.5}{min}），从而无谓消耗续航并把本可保障的架次误判为不可行。故中继须在\textbf{每个区间起点之前}完成建链，服务能耗按各区间被覆盖时长之和计算，重叠区间先合并以免重复计费。

\subsection{中继排班：把服务窗口作为硬约束}

\textbf{模型}：给定每个需保障架次的中断区间集合，为其分配（悬停点、中继无人机、能源组件、起飞时刻），满足：

\begin{enumerate}[label=(\arabic*),leftmargin=2em]
	\item \textbf{窗口硬约束}：中继必须在其负责的\textbf{每个}区间起点之前完成建链。资源来不及 ⇒ 该架次记为\textbf{不可行}，\textbf{绝不放宽窗口}；
	\item \textbf{共享（A 口径）}：允许\textbf{同一悬停点被多架运输机共享}（依据见 \S\ref{sec:share}）；
	\item \textbf{周转}：同一中继无人机相邻架次之间须留出架次周转时间；能源组件须\textbf{充电完成}后才能再次投入。
\end{enumerate}

求解策略为按区间起点排序的\textbf{贪心}：对每个需求，优先尝试并入已有中继架次（同点共享，按区间\textbf{并集}评估），否则新开一个架次并选“能赶上且能耗最低”的 (悬停点, 无人机, 能源组件) 组合。

\subsection{中继共享口径的依据}\label{sec:share}

本文允许同一悬停点被多架运输机共享，依据是\textbf{题目附录原文}：

\begin{quote}
	每架\textbf{运输无人机}在任一时刻只能由 G01 或\textbf{一架}中继无人机提供通信保障。
\end{quote}

这是对\textbf{运输机侧}的约束（每架运输机同时刻至多一个保障源），\textbf{并未}限制一架中继的服务对象数量。因此“同点共享”是题目允许的，而非对约束的放宽。

\subsection{数值结果（本文结果）}

\subsubsection{直连诊断：中继是否必需}

沿每条架次的完整轨迹逐时刻采样并判定直连可用性，实测 25 个运输架次中 \textbf{20 个存在直连中断}（占 80.0\%、平均中断占比 \textbf{28.4\%}），另 5 个架次全程直连可用。这直接证明\textbf{中继无人机是必需项而非可选项}——若没有中继，这 20 个架次在中断时段将失去指挥与遥测链路。

\subsubsection{覆盖率：两个不可混用的口径}\label{sec:q3gap}

这是本问\textbf{最容易失分}的地方，必须严格区分：

\begin{table}[htbp]
	\caption{问题三覆盖率与资源缺口（本文结果）}
	\label{tab:q3cov}
	\centering
	\zihao{-5}
	\begin{tabular}{lc}
		\toprule
		指标 & 数值 \\
		\midrule
		运输架次 / 其中需保障       & 25 / \textbf{20} \\
		中继架次                    & \textbf{6} \\
		\textbf{几何可达覆盖率}     & $20/20=100\%$（\textbf{仅选址能力上界}） \\
		\textbf{时间轴覆盖率}       & $\mathbf{11/20=55.0\%}$（\textbf{唯一可上报口径}） \\
		中继资源缺口                & \textbf{4 架}（需 6 个中继架次 vs 库存 2 架） \\
		运输 / 中继 / 总能耗        & 77.30 / \textbf{2.35} / \textbf{79.65}\,kWh \\
		联合任务完成时间            & \textbf{10968.7}\,s（3.05\,h） \\
		独立校验器                  & 硬违规 \textbf{0 条}；软违规 42 条（均为中继资源不足） \\
		\bottomrule
	\end{tabular}
\end{table}

\textbf{几何可达覆盖率}只回答“\textbf{是否存在}一个悬停点能覆盖该架次全部中断样本”，是选址能力的\textbf{上界}；\textbf{时间轴覆盖率}才回答“中继\textbf{那一刻是否在站}”。本文以\textbf{时间轴口径}为唯一上报依据。

\textbf{缺口的定量根因。}缺口来自\textbf{题目给定的资源配置}，不是求解算法的缺陷。由式~\eqref{eq:hovermax}，单架中继的\textbf{在站时长上限约 \SI{140}{min}}（尚未扣除往返飞行能耗）；而 20 个需保障架次的\textbf{真实中断时长合计约 \SI{168.5}{min}}，\textbf{最大并发达 6 个架次}。仅 2 架中继在时间轴上无法同时满足。因此覆盖率必然低于几何可达覆盖率，本文如实报告 \textbf{55.0\%} 与 \textbf{4 架中继缺口}。

\textbf{改进方向}有两条：\textbf{其一}，将中继无人机数量由 2 架增至 6 架即可消除缺口；\textbf{其二}，在现有资源下做\textbf{时空联合优化}，把有限的在站时长优先投给“单位在站时间可覆盖架次数最多”的悬停点，以最大化获保障架次数。本文采用后者作为求解策略，并如实报告剩余缺口。

\subsubsection{选址规律}

所有中继悬停点的离地高度均取上限 \SI{250}{m}（离地越高视线越好），悬停海拔介于 \SIrange{440}{962}{m}；水平位置集中在服务区群中心偏西（约 $109.17^\circ\sim109.27^\circ$E），即\textbf{贴近作业空域而非贴近网关}。这与链路门限分析完全一致：三条链路门限为 直连 \SI{122}{dB}、中继接入 \SI{116}{dB}、中继回传 \SI{126}{dB}，其中\textbf{接入段最弱}（比直连低 \SI{6}{dB}），故中继必须靠近运输机；而回传段宽松（可达 \SI{19.83}{km}），对位置不敏感。

\subsection{一次被纠正的严重误报（方法论记录）}\label{sec:q3bug}

本文早期版本曾报告“中继覆盖 $20/20=100\%$”。该结论\textbf{是错误的}，且错误被\textbf{静默吞掉}，根因有三处叠加：

\begin{enumerate}[label=(\arabic*),leftmargin=2em]
	\item \textbf{服务区间被静默截短}：评估函数取 $t_{svc}^{start}=\max(t_{link},T_{win}^{start})$，中继晚到时只会把服务区间\textbf{截短}（截到 0 长度也不报错），于是“中继在运输机返航之后才到场”被当成已保障；
	\item \textbf{排班未把服务窗口纳入约束}：只按“最早可用资源”定起点，波次 1 有 8 个运输架次同时起飞而中继只有 2 架，窗口必然冲突；
	\item \textbf{校验器从未执行该项检查}（这是缺陷被吞掉的\textbf{直接原因}）：旧实现对“已安排中继”的架次\textbf{不填中断窗口字段}，于是校验器直接跳过通信检查，报了 0 违规。
\end{enumerate}

\textbf{修正}：① 评估函数改为\textbf{显式失败}（无法在区间起点前建链即判定不可行）；② 排班把服务窗口作为\textbf{硬约束}；③ 一律如实上报真实中断区间，使校验器的通信检查\textbf{真正生效}。修正后时间轴覆盖率由 $0/20$ 变为 $11/20$。

\textbf{方法论教训}：\textbf{“校验器报 0 违规”不等于“方案没有问题”}——必须先确认校验器\textbf{确实执行了该项检查}。本文因此把校验器改为区分两类违规：\textbf{资源缺口}（软违规，如实报告、不阻断产出）与\textbf{排班缺陷}（硬违规，闸门拦截）。若把资源缺口也当硬违规，求解器会直接中止、什么都不产出，等于“因为资源不够就干脆不报告”。

\textbf{本问适用条件}：结论依赖 2 架中继的库存与其中继参数（式~\eqref{eq:hovermax}）；若中继数量改变，覆盖率需重算。$\Delta t=\SI{2}{s}$ 的采样步长经敏感性分析确认足以捕捉遮挡特征（由 \SI{5}{s} 细化到 \SI{2}{s} 时中断占比判定变化小于 2\%）。
